% ============================================================
% RESUMEN
% ============================================================

\begin{abstract}

Este informe documenta el diseño e implementación de un sistema de digitalización de señales basado en un microcontrolador de 32 bits (ESP32-WROOM, vía el núcleo Arduino para ESP32; la guía de práctica nombra una tarjeta STM32 específicamente para la etapa de reconstrucción PWM, pero este proyecto se desarrolla sobre ESP32 al ser la tarjeta realmente disponible — ver Sección de Metodología para las diferencias de hardware relevantes, en particular la resolución del DAC integrado), desarrollado para la práctica \texttt{LAB1} del curso. El sistema abarca cuatro etapas: (1) digitalización de señales analógicas periódicas mediante interrupción por temporizador a $500$ muestras/segundo con conversión análogo-digital de $12$ bits; (2) reconstrucción de la señal digitalizada mediante el módulo DAC de la tarjeta y, alternativamente, mediante PWM a $5$ kHz seguido de un filtro pasa-bajos analógico externo; (3) digitalización de señales de sensores —un sensor de presión, temperatura y humedad BME280 por I2C, a $100$ muestras/segundo (límite impuesto por el propio sensor, no $200$ como en la sección de encoder — ver Sección de Metodología), y un encoder acoplado a un motorreductor DC a $200$ muestras/segundo—; y (4) identificación experimental del fenómeno de aliasing con frecuencias superiores a la de muestreo. La guía de práctica nombra originalmente una unidad de medición inercial (IMU) para esta tercera etapa; el sensor realmente disponible es un BME280 en su lugar, por lo que esta etapa mide presión/temperatura/humedad y no movimiento — ver Sección de Metodología para el detalle. Todo el firmware, el programa de captura eficiente hacia el computador (sin visualización en tiempo real, por requisito de la práctica) y los scripts de análisis en MATLAB (FFT, estadística descriptiva, generación de las tablas solicitadas) están implementados y documentados en este repositorio. \textbf{No se dispuso de hardware físico} (tarjeta ESP32-WROOM, osciloscopio, generador de señales, BME280 ni motorreductor) en el entorno de desarrollo de este informe; en consecuencia, las tablas de resultados se presentan con sus valores marcados explícitamente como pendientes de captura real, y las preguntas de reflexión de la guía se responden desde el fundamento teórico del muestreo, la cuantización y la reconstrucción de señales, en lugar de a partir de mediciones inventadas. Todo el código de este proyecto está disponible en el repositorio público en \url{https://github.com/DanielAraqueStudios/Signals-frecuency/tree/main/LAB_TESTING/LAB1}.

\end{abstract}

% ============================================================
% PALABRAS CLAVE
% ============================================================

\begin{IEEEkeywords}

Muestreo, cuantización, aliasing, conversión análogo-digital, conversión digital-análoga, PWM, filtro pasa-bajos, BME280, encoder, ESP32, Transformada Rápida de Fourier.

\end{IEEEkeywords}
